\chapter*{Abstract}
\addcontentsline{toc}{chapter}{Abstract}

\glsunset{ADNEX}
\glsunset{RMI}
\glsunset{AUC}

Clinical prediction models can support evidence-informed clinical decision-making by estimating a patient's risk of disease. Their safe and effective use, however, requires more than good discrimination on average: predicted probabilities must be accurate (i.e. calibrated), clinically useful, and sufficiently robust to capture the differences between centres, measurement procedures, and patient populations.


This thesis combines applied and methodological research to improve decision-making for women with adnexal masses, with a focus on the Assessment of Different NEoplasias in the adneXa (\gls{ADNEX}) model for ovarian cancer diagnosis.

First, we synthesised the external validation evidence for \gls{ADNEX}. A systematic review and meta-analysis of 47 studies including more than 17\,000 tumours confirmed strong discrimination (AUC 0.93 for benign versus malignant masses), but also showed that the methodological reporting quality is poor: 91\% of validations were at high risk of bias and calibration was rarely assessed. Second, in a head-to-head systematic review and meta-analysis of 11 studies (8271 tumours), \gls{ADNEX} consistently outperformed the Risk of Malignancy Index (\gls{RMI}) in discrimination and clinical utility; at a 10\% malignancy-risk threshold, the probability of being useful in a hypothetical new centre was 96\% for \gls{ADNEX} and 15\% for \gls{RMI}. Third, we updated \gls{ADNEX} to \emph{ADNEX2} using multicentre prospective data that also included conservatively managed patients, better adapting the model's use in a two-step strategy that starts with the benign simple descriptors.

The methodological part of this thesis addresses broader challenges in probability estimation and model evaluation. We studied why random forest models with excellent training performance can still yield unstable or misleading probabilities, and showed how random forests can learn local probability peaks for which common defaults (fully grown trees) may be counterproductive when the goal is well-calibrated risk estimation. We then developed clustered calibration methods that explicitly incorporate multicentre structure, allowing calibration to be summarised not only on average but also with prediction intervals that quantify expected between-cluster variability. Next, we proposed a framework for uncertainty in individual risk estimates and illustrated empirically that, even with large training samples, the uncertainty is uncomfortably large. This supports the idea that models should always be interpreted with caution and that decisions taken by expert clinicians can place the risk estimates into context. Finally, we introduced a variable selection algorithm that optimizes clinical utility, ensuring that the selected predictors are clinically relevant while taking into account the cost of measuring them.

The thesis advances the field of clinical prediction modelling by providing new evidence on the performance of \gls{ADNEX} and creating \emph{ADNEX2}, and by developing novel methods for prediction model development and validation that can be applied in a wide range of clinical contexts. The findings have important implications to guide the development and validation of clinical prediction models towards trustworthy and useful tools for clinical decision-making.